\section{Текущие результаты}

Реализованный пайплайн уже поддерживает сбор корпуса, генерацию embeddings, извлечение ключевых слов, построение prompts, retrieval референсов, генерацию кандидатов и pairwise automatic evaluation. Оставшийся отсутствующий этап --- собственно reinforcement-learning update loop. Поэтому данный раздел является узким: он описывает завершенные data artifacts и фиксирует конкретный experimental setting для первых MRVF training runs.

Результат на уровне корпуса --- объединенная коллекция шуток примерно из \(664\text{k}\) объектов. В финальной training setup этот полный корпус служит retrieval pool: все шутки кодируются embeddings и индексируются, а references для weak prompts извлекаются из этого полного embedding space. Однако keyword extraction нужно применять только к достаточно большому подмножеству, чтобы достичь целевого sample size. В нашем случае \(3-5\text{k}\) keyword groups могут дать \(10-20\text{k}\) references для training, validation и testing.

Для preliminary stage этого достаточно. Важный результат состоит в том, что данные и experimental setting теперь финализированы: корпус существует, reference collection реализована, методология зафиксирована, а оставшаяся работа ограничена MRVF training loop.

\section{Планируемая процедура обучения}

Репозиторий уже реализует полный data pipeline на Python с использованием Hugging Face Datasets, parquet artifacts, FAISS-based nearest-neighbor search и OpenAI-compatible API calls для embedding generation, candidate generation и automatic evaluation. Однако для собственно training runs software stack немного изменится.

Планируемая реализация будет использовать patched version of \texttt{trl} для optimization и \texttt{vLLM} для generation. \texttt{trl} предоставляет trainer-side logic для policy optimization, тогда как \texttt{vLLM} будет использоваться и для генерации baseline outputs, и для обслуживания rollouts во время reinforcement learning. Поскольку MRVF зависит от sampling нескольких reasoning traces на prompt, rollout throughput важен для эффективности.

Целевые модели относятся к семейству Qwen3 в диапазоне размеров \(1.7\text{B}\) и \(4\text{B}\). Для каждого размера планируются следующие baselines:
\begin{enumerate}[leftmargin=1.8em]
    \item \textbf{Qwen3 Instruct}, baseline в non-thinking mode;
    \item \textbf{Qwen3 Thinking}, baseline в thinking mode;
    \item \textbf{MRVF-trained Qwen3 Thinking}, предлагаемая модель;
\end{enumerate}

Такой сеттинг сохраняет интерпретируемость сравнения. Он изолирует три отдельных вопроса: достаточно ли обычного instruction tuning для генерации юмора; улучшает ли native thinking mode генерацию юмора уже без дополнительного обучения; и может ли MRVF дополнительно улучшить thinking model.

Все training runs планируются на 1--2 \texttt{NVIDIA H200} GPUs в HSE HPC cluster. Этого достаточно для целевых размеров моделей и позволяет распараллеливать эксперименты между instances.

Сравнение будет использовать evaluation pipeline, уже реализованный в репозитории. Outputs сопоставляются для одного prompt, left/right order рандомизируется, judge принуждается выбрать ровно одного winner, а model scores агрегируются через Bradley-Terry leaderboard. Этого достаточно для эксперимента, поскольку цель состоит в обнаружении того, дает ли MRVF измеримый сдвиг относительно instruct и thinking baselines. Более крупная multidimensional evaluation с rubric-based scoring и human annotation планируется после завершения обучения.

\section{Планируемые ablations}

Ablation plan напрямую связан с теорией. Наиболее информативные ablations следующие:
\begin{itemize}[leftmargin=1.5em]
    \item \textbf{Число наблюдаемых референсов на prompt \(|Y^{*}|\).} Сравнить меньшие и большие reference sets, например через truncating retrieval до \(1\), \(3\) или \(5\) объектов.
    \item \textbf{KL coefficient \(\beta\).} Проверить, насколько сильно модель должна быть anchored to initialization checkpoint, чтобы компенсировать incompleteness.
    \item \textbf{Group size \(G\).} Варьировать число sampled reasoning traces per prompt, чтобы оценить variance / throughput trade-off MRVF estimator.
    \item \textbf{Model size.} Сравнить \(1.7\text{B}\) и \(4\text{B}\), а также проверить, насколько MRVF устойчив при применении к меньшим моделям.
\end{itemize}
